\subsection{円錐基礎枝の排除と標準AS3部分代数}
\label{subsec:h3-positive3-cone-exclusion}
\begin{lemma}\label{lem:h3-curve-graded-division-growth}
$K=k(u)$、$\sigma\in\operatorname{Aut}_kK$とする。
次数1の有限個の斉次元で生成される
$K[h,h^{-1};\sigma]$の部分代数はGK次元2以下である。
従って円錐基礎$S/(C')$の場合、残る正生成元$v$が共通量子平面$T=\langle W\rangle$の
斉次商体に入るH3包絡は存在しない。
\end{lemma}
\begin{proof}
各生成元を$f_i(u)h$と書き、その全極を支配する有効因子$D$を$\mathbb P^1$上に取る。
長さ$n$の語の係数の極は$D_n=\sum_{j=0}^{n-1}\sigma^jD$で抑えられる。
$\deg D_n=n\deg D$なので、係数空間の次元は$O(n)$、累積成長は$O(n^2)$である。
円錐基礎の斉次商環は、前の量子torus表示から
$k(u)[z,z^{-1};u\mapsto au]$に等しい。
もし$v\in Q_{\mathrm{gr}}(T)$なら包絡の四生成元が全てこの斉次商環に入り、
14641のGK次元4に反する。
\end{proof}

\begin{theorem}\label{thm:h3-positive3-cone-base-exclusion}
定理\ref{thm:h3-positive3-canonical-base-coverage}の量子Veronese円錐候補は空である。
従って全$(3,1)$H3対象の標準基礎$A_0=\langle W,z\rangle$は、
対角Heisenberg型またはJordan型のquadratic AS3として$B$と$R$の双方へ埋まる。
\end{theorem}
\begin{proof}
円錐の場合と仮定する。その関係は
\[
 T=(xy=a^{-2}yx),\quad zx=axz,\quad zy=a^{-1}yz,
 \quad z^2=\kappa xy,\quad\kappa\ne0,\quad a^4\ne1.
\]
$W$上の中山は恒等で、残る正生成元は$v$である。
二つの追加二次関係を$v^2$へ射影する。
像が零なら、両cross射影が可逆でない場合に$v\ell=q(W)$または$\ell v=q(W)$となり、
前補題で排除される。
両射影が可逆なら$v$は$T$を自己同型$g$とskew derivationでOre正規化する。
その互換性は$T$の関係に$v$を作用させ、前補題で$v\notin Q_{\mathrm{gr}}(T)$を使えば従う。
四つの基礎関係と二つのOre cross式には、外側の$vz$や$v^2$を減らす式がない。
$T$のordered monomialと、$zz$を避ける$v,z$語を用いる書換えは全overlapを解消する。
必要な条件は$f^2(t)\kappa xy=\kappa xy\,t$、$f(xy)=xy$と$g$-derivation条件だけである。
従ってそのHilbert級数は
\[
 \frac1{(1-t)^2}\frac{1+t}{1-t-t^2}
\]
となり、次数3の値21が14641値20に反する。

よって$v^2$射影は一次元である。
二つの関係のうち一つは$v^2$係数1、もう一つはcrossと$W^2$だけを持つ。
後者の片方のcross射影が零なら再び$v$が商体へ落ちるため、両方非零である。
四次topの$v^4$係数は中山巡回式から零で、$v^2$行は非零である。
もし$v^3W$成分が零なら、この行は$W^2$内の唯一の関係$r_T$の非零倍数となる。
一方、全cross行は$v^2,z^2$係数が零なので、平方射影の核の二本
（$r_T$とcross関係）に入り、cross--cross blockの階数は高々1。
正三元のanti-cyclic式は、$v^2$--$W^2$ blockがそのcross非対角blockの逆符号であることを与える。
従って前者の階数も高々1となるが、$r_T$の二変数係数行列は階数2なので矛盾する。
ここで中山に$v\mapsto v+w$という単冪成分があっても、
$v^3W=0$の下では二つの$v$を持つ成分にその補正は現れず、同じanti-cyclic式が成り立つ。

特に正三次元中山拡大は分裂する。
非分裂なら$Dv=w\in W,DW=0$であり、既出の不変tensor恒等式により$v^3W$が零となって、
直前の矛盾を生じるからである。
従って中山は半単純であり、$v^3W\ne0$である。
$v^2$行のcross部分を$n(v\ell-\ell v)$、$n\ne0$と書く。

偶関係の基底を
\[
 r_T,\quad z^2-\kappa xy,\quad
 r_{\mathrm{cross}},\quad r_{v^2}
\]
と取り、最後の関係だけが$v^2$係数1を持つようにする。
前二本のcross部分は零、後二本のcross部分を$r=(\ell,-\ell),s$とする。
middle Frobenius行列をこの基底で$M$と書くと、
$v^4$係数が零だから$M_{44}=0$である。
従ってcross--cross blockは
\[
 H_{\mathrm{cross}}=M_{33}rr^t+M_{34}(rs^t+sr^t)=rk^t+kr^t
\]
と書ける。
これは定理\ref{thm:h3-positive3-global-cross}と同じanti-cyclic blockであり、
両対角blockの逆符号から$k$の両成分が一致する。
$v$の$W$平行移動で全cross blockと$v^2$--$W^2$ blockを消せるため、
\[
 [v,\ell]=0,\qquad v^2\in T_2
\]
となる。$z^2=\kappa xy$は$v$を含まないので、この基底変更でも変わらない。
$a^4\ne1$は$T$の量子パラメータ$a^{-2}$が$\pm1$でないことを意味する。
従って$\ell$の二次centralizerは$k\ell^2$だけで、$v^2=c\ell^2$。
可換する$v,\ell$の平方差が因子化して整域性に反する。
これで円錐枝は排除された。
\end{proof}

\begin{corollary}\label{cor:h3-positive3-AS3-base-Hilbert}
全$(3,1)$H3対象について標準AS3部分代数は$R$にも埋まり、
\[
 \frac{H_R(t)}{H_{A_0}(t)}=1+t+t^2.
\]
これはHilbert級数の等式であり、$R$が$A_0$上有限自由階数3であることを
この等式だけから結論しない。
\end{corollary}
\begin{proof}
基礎はquadratic AS3なので$H_{A_0}=(1-t)^{-3}$、
$H_R=(1-t^3)(1-t)^{-4}$である。
\end{proof}

\texttt{verify\_cubic\_positive3\_cone\_exclusion.py}は
次数3の成長値、Frobenius関係基底からのblock式、およびcentralizer minor
$-(a^4-1)^2$を厳密検算する。
用いたanti-cyclic基底変更と非分裂tensorについても、既存の全係数検証器を再実行した。
このAS3基礎の存在から、追加正生成元の全関係やH3許容性が済んだとはしない。
